% Part: sets-functions-relations
% Chapter: relations
% Section: orders

\documentclass[../../../include/open-logic-section]{subfiles}

\begin{document}

\olfileid{sfr}{rel}{ord}
\olsection{序}

\begin{explain}
我们的许多比较都涉及在某个方面将一些对象描述为“小于”、“等于”或“大于”另一些对象。这些比较涉及\emph{序}关系。但序关系有多种不同的类型。例如，有些要求任何两个对象都是可比的，有些则不要求；有些包含等同关系（如~$\le$），有些则排除等同关系（如~$<$）。在此建立一个分类法将会对我们有所帮助。
\end{explain}

\begin{defn}[预序]
一个既自反又传递的关系称为\emph{预序}。  
\end{defn}

\begin{defn}[偏序]
一个同时也是反对称的预序称为\emph{偏序}。
\end{defn}

\begin{defn}[全序]\ollabel{def:linearorder}
一个同时也是连通的偏序称为\emph{全序}或\emph{线性序}。
\end{defn}

\begin{ex}
每个线性序都是偏序，每个偏序也都是预序，但反之则不成立。集合~$A$ 上的全集关系是一个预序，因为它是自反且传递的。但是，如果 $A$ 的元素多于一个，那么全集关系就不是反对称的，因而也不是偏序。
\end{ex}

\begin{ex}
考虑~$\Bin^*$ 上的\emph{不长于}关系 $\preccurlyeq$：$x \preccurlyeq y$ 当且仅当 $\len{x} \le \len{y}$。这是一个预序（自反且传递），甚至是连通的，但不是偏序，因为它不是反对称的。例如，$01 \preccurlyeq 10$ 且 $10 \preccurlyeq 01$，但是 $01 \neq 10$。
\end{ex}

\begin{ex}
一个重要的偏序是集合族上的包含关系 $\subseteq$。它一般不是线性序，因为如果 $a \neq b$，我们考虑 $\Pow{\{a, b\}} = \{\emptyset, \{a\}, \{b\}, \{a,b\}\}$，
就会看到 $\{a\} \nsubseteq \{b\}$ 且 $\{a\} \neq \{b\}$ 以及 $\{b\}
\nsubseteq \{a\}$。
\end{ex}

\begin{ex}
\emph{无余数整除}关系给出了一个不是线性序的偏序。对于整数 $n$ 和~$m$，我们写作 $n \mid m$ 表示 $n$（整）除 $m$，即当且仅当存在某个整数~$k$ 使得 $m = kn$。在~$\Nat$ 上，这是一个偏序，但不是线性序：例如，$2 \nmid 3$ 且 $3 \nmid 2$。若将整除视为 $\Int$ 上的关系，则它只是一个预序，因为它不是反对称的：$1 \mid -1$ 且 $-1 \mid 1$，但 $1 \neq -1$。
\end{ex}

\begin{ex}
序列集~$A^*$ 上的\emph{扩展}关系定义如下：$s \sqsubseteq s'$ 当且仅当 $s = \emptyseq$（空序列），$s = s'$，或者 $s = \tuple{s_1, \dots, s_n}$ 且 $s' =
\tuple{s_1, \dots, s_n, s_{n+1}, \dots, s_m}$。如果 $s \sqsubseteq s'$，我们也说 $s$ 是~$s'$ 的一个\emph{初始段}。$A^*$ 上的扩展关系是偏序但不是线性序，例如，如果 $a \neq b$，那么 $ab \not\sqsubseteq ba$ 且 $ba \not\sqsubseteq ab$。
\end{ex}

\begin{defn}[严格序]
\emph{严格序}是一个反自反的、非对称的且传递的关系。
\end{defn}

\begin{defn}[严格线性序]\ollabel{def:strictlinearorder}
一个同时也是连通的严格序称为\emph{严格全序}或\emph{严格线性序}。
\end{defn}

\begin{ex}
$\le$ 是对应于严格线性序~$<$ 的线性序。$\subseteq$ 是对应于严格序~$\subsetneq$ 的偏序。
\end{ex}

任何~$A$ 上的严格序 $R$ 都可以通过添加对角线 $\Id{A}$，即添加所有序对~$\tuple{x, x}$，而变成一个偏序。（这称为 $R$ 的\emph{自反闭包}。）反之，从一个偏序出发，通过移除~$\Id{A}$ 可以得到一个严格序。下面两个结果精确地表述了这一点。

\begin{prop}\ollabel{prop:stricttopartial}
如果 $R$ 是~$A$ 上的一个严格序，那么 $R^+ = R \cup \Id{A}$ 是一个偏序。此外，如果 $R$ 是一个严格线性序，那么 $R^+$ 是一个线性序。
\end{prop}

\begin{proof}
假设 $R$ 是一个严格序，即 $R \subseteq A^2$ 且 $R$ 是反自反的、非对称的、传递的。令 $R^+ = R \cup \Id{A}$。我们需要证明 $R^+$ 是自反的、反对称的和传递的。

$R^+$ 显然是自反的，因为对所有 $x \in A$ 都有 $\tuple{x, x} \in \Id{A} \subseteq
R^+$。 

为证 $R^+$ 反对称，反设 $R^+xy$ 且 $R^+yx$ 但 $x \neq y$。由于 $\tuple{x,y} \in R \cup \Id{A}$，而 $\tuple{x, y} \notin \Id{A}$，我们必有 $\tuple{x, y} \in R$，即 $Rxy$。类似地，有 $Ryx$。但这与 $R$ 的非对称性矛盾。

要证传递性，假设 $R^+xy$ 且 $R^+yz$。如果 $\tuple{x, y} \in R$ 且 $\tuple{y,z} \in R$，那么由 $R$ 的传递性得 $\tuple{x, z} \in R$。否则，要么 $\tuple{x, y} \in
\Id{A}$，即 $x = y$，要么 $\tuple{y, z} \in \Id{A}$，即 $y = z$。在第一种情况下，由假设知 $R^+yz$ 且 $x = y$，故 $R^+xz$。第二种情况同理。无论哪种情况，都有 $R^+xz$，故 $R^+$ 也是传递的。

关于“此外”的结论，假设 $R$ 也是连通的。那么对所有 $x \neq y$，要么 $Rxy$，要么 $Ryx$，即要么 $\tuple{x, y} \in R$，要么 $\tuple{y, x} \in R$。由于 $R \subseteq R^+$，这对 $R^+$ 依然成立，故 $R^+$ 也是连通的。
\end{proof}

\begin{prop}\ollabel{prop:partialtostrict}
如果 $R$ 是 $A$ 上的一个偏序，那么 $R^- = R \setminus \Id{A}$ 是一个严格序。此外，如果 $R$ 是一个线性序，那么 $R^-$ 是一个严格线性序。
\end{prop}

\begin{proof}
留作练习。
\end{proof}

\begin{prob}
请证明 \olref[sfr][rel][ord]{prop:partialtostrict}。 
\end{prob}

下述简单结果表明，严格线性序满足一种类似外延性的性质：

\begin{prop}\ollabel{prop:extensionality-strictlinearorders}
如果 $<$ 是 $A$ 上的一个严格线性序，那么：
\[
  (\forall a, b \in A)((\forall x \in A)(x < a \liff x < b) \lif a = b).
\]
\end{prop}

\begin{proof}
假设 $(\forall x \in A)(x < a \liff x < b)$。如果 $a < b$，那么 $a < a$，这与 $<$ 的反自反性矛盾；所以 $a \nless b$。同理可得 $b \nless a$。又因为 $<$ 是连通的，所以 $a = b$。
\end{proof}

\end{document}